\clearpage
\section{Аналитическая часть}

\subsection{Постановка задачи}

Дан неориентированный граф $G=(V,E)$ с $n=|V|\geq1$. Вершины соответствуют портам, рёбра --- возможным морским переходам. Каждому ребру $\{i,j\}$ сопоставлены положительное целое время $t_{ij}=t_{ji}$ в часах и положительная целая стоимость $c_{ij}=c_{ji}$ в условных денежных единицах. Лимит на весь маршрут обозначен через $B\geq0$.

Открытый маршрут $P=(v_1,\ldots,v_n)$ представляет собой перестановку вершин графа. Начальный и конечный порты свободны. Для всех $k=1,\ldots,n-1$ требуется $\{v_k,v_{k+1}\}\in E$. Целевая функция и бюджетное ограничение имеют вид
\begin{equation}
\begin{aligned}
L(P)&=\sum_{k=1}^{n-1}t_{v_kv_{k+1}}\longrightarrow\min,\\
C(P)&=\sum_{k=1}^{n-1}c_{v_kv_{k+1}}\leq B.
\end{aligned}
\label{eq:objective}
\end{equation}
При одинаковом времени предпочтение отдаётся меньшей стоимости. В обратном порядке тот же набор рёбер имеет те же показатели. При $n=1$ единственная вершина образует допустимый маршрут с нулевыми затратами и временем. Экспериментальные графы полные. Программа также поддерживает графы с отсутствующими переходами.

\subsection{Метод полного перебора}

Перестановки портов генерируются в лексикографическом порядке. Для каждой суммируются время и стоимость переходов. Проверка перестановки прекращается при отсутствии ребра или превышении бюджета.

Преимущество метода --- гарантия оптимальности найденного маршрута и возможность доказать отсутствие решения. Недостаток --- факториальный рост числа вариантов. Для $n$ портов перебирается $n!$ перестановок, включая развороты маршрутов.

\clearpage
\subsection{Муравьиный метод с элитным усилением}

Каждый муравей случайно выбирает начальный порт и строит маршрут по одному переходу. Для каждого ребра программа хранит значение феромона. Это число повышает вес перехода при положительном влиянии феромона и прочих равных условиях. На выбор также влияют время перехода и стоимость относительно оставшегося бюджета. После каждой итерации значение на ребре уменьшается на долю испарения, затем к нему прибавляется вклад допустимых маршрутов. Муравьи, построившие допустимые маршруты, вносят феромон на их рёбра. Элитные муравьи дополнительно усиливают рёбра лучшего допустимого маршрута, найденного с начала поиска. Этот принцип усиления предложен в работе Дориго, Маньеццо и Колорни~\cite{ant}.

Преимущество метода состоит в возможности ограничить время поиска числом муравьёв и итераций. Накопленный феромон повышает вероятность повторного использования удачных переходов. Недостаток метода --- случайный характер результата. За отведённое число попыток алгоритм может получить маршрут с большим временем или пропустить все допустимые маршруты. В частности, это возможно при жёстком бюджете.

\subsection*{Вывод}

Полный перебор служит источником эталонного времени, а метод с элитным усилением позволяет исследовать качество маршрутов при ограниченном числе попыток. Оба метода используют общую целевую функцию и одно бюджетное ограничение.
